% ===== PRÉAMBULE CANONIQUE — à coller VERBATIM en tête de CHAQUE .tex =====
\documentclass[11pt,a4paper]{article}
\usepackage[utf8]{inputenc}\usepackage[T1]{fontenc}\usepackage{lmodern}
\usepackage[french]{babel}
\usepackage{amsmath,amssymb,mathtools}\usepackage{icomma}
\usepackage{siunitx}\sisetup{locale=FR}  % \qty{}{}, \si{}, \ang{}
\usepackage{xcolor}\usepackage{colortbl}
\usepackage{tikz}\usepackage{pgfplots}\pgfplotsset{compat=1.18}
\usepackage[siunitx,europeanresistors]{circuitikz} % circuits (élec.)
\usepackage[most]{tcolorbox}\usepackage{enumitem}\usepackage{array,booktabs}
\usepackage[a4paper,margin=2.2cm]{geometry}
\usetikzlibrary{arrows.meta,calc,angles,quotes,positioning,patterns,%
  backgrounds,shapes.geometric,decorations.pathmorphing,decorations.markings,intersections}
\definecolor{primary}{RGB}{222,49,99}\definecolor{primarydark}{RGB}{150,22,55}
\definecolor{defcol}{RGB}{0,114,178}\definecolor{methcol}{RGB}{52,92,148}
\definecolor{excol}{RGB}{0,135,90}\definecolor{attncol}{RGB}{200,40,55}
\tcbset{boxbase/.style={enhanced,breakable,boxrule=0.9pt,arc=2.5pt,
  left=3mm,right=3mm,top=2.3mm,bottom=2.3mm,fonttitle=\bfseries,coltitle=white}}
% --- boîtes TUTORIEL (noms EXACTS, ne pas renommer) ---
\newtcolorbox{cle}[1][]{boxbase,colback=defcol!6,colframe=defcol,
  title={\textsc{À retenir}\ifx\relax#1\relax\relax\else\textmd{ : \itshape #1}\fi}}
\newtcolorbox{methode}[1][]{boxbase,colback=methcol!7,colframe=methcol,sharp corners,
  title={\textsc{Méthode}\ifx\relax#1\relax\relax\else\textmd{ --- \itshape #1}\fi}}
\newtcolorbox{exemplebox}[1][]{boxbase,colback=excol!5,colframe=excol,
  title={\textsc{Exemple}\ifx\relax#1\relax\relax\else\textmd{ : \itshape #1}\fi}}
\newtcolorbox{piege}[1][]{boxbase,colback=attncol!6,colframe=attncol,
  title={\textsc{Piège}\ifx\relax#1\relax\relax\else\textmd{ : \itshape #1}\fi}}
\newtcolorbox{remarque}[1][]{boxbase,colback=black!4,colframe=black!45,
  title={\textsc{Remarque}\ifx\relax#1\relax\relax\else\textmd{ : \itshape #1}\fi}}
% --- boîtes EXERCICE / CORRIGÉ (noms EXACTS, auto-comptées) ---
\newtcolorbox[auto counter]{exo}[1][]{boxbase,colback=primary!4,colframe=primary,
  title={\textsc{Exercice}~\thetcbcounter\ifx\relax#1\relax\relax\else\textmd{ --- \itshape #1}\fi}}
\newtcolorbox[auto counter]{corrige}[1][]{boxbase,colback=excol!4,colframe=excol,
  title={\textsc{Corrigé}~\thetcbcounter\ifx\relax#1\relax\relax\else\textmd{ --- \itshape #1}\fi}}
\newcommand{\vv}[1]{\vec{#1}}\newcommand{\ds}{\displaystyle}
\newcommand{\R}{\mathbb{R}}\newcommand{\N}{\mathbb{N}}\newcommand{\Z}{\mathbb{Z}}
% (un \usepackage{fancyhdr}+en-tête est possible ; il est ignoré côté web)
% =========================================================================
\begin{document}
\begin{exo}[calculer la force de Laplace]
Un conducteur rectiligne de longueur $L=\qty{15}{\centi\metre}$ est parcouru par un courant
$I=\qty{3}{\ampere}$. Il est placé \textbf{perpendiculairement} à un champ magnétique uniforme
$B=\qty{0.4}{\tesla}$.
\begin{center}
\begin{tikzpicture}[>=Stealth,scale=1.0]
  \foreach \y in {-1,-0.6,-0.2,0.2,0.6,1}{\draw[blue!60!black,line width=0.8pt,->] (-2.4,\y)--(-1.6,\y);}
  \node[blue!70!black] at (-2.7,0){$\vv{B}$};
  \draw[red!80!black,line width=2.2pt] (-0.4,-1.1)--(-0.4,1.1);
  \draw[red!80!black,->,line width=2.2pt] (-0.4,0.1)--(-0.4,0.7); \node[red!80!black,right] at (-0.3,0.9){$I$};
  \node[font=\footnotesize] at (0.9,0){$L,\ \theta=\ang{90}$};
\end{tikzpicture}
\end{center}
Calcule l'intensité de la force de Laplace subie par le conducteur.
\end{exo}

\begin{exo}[le conducteur parallèle au champ]
Un conducteur rectiligne parcouru par un courant $I=\qty{6}{\ampere}$ est placé
\textbf{parallèlement} aux lignes d'un champ $B=\qty{0.5}{\tesla}$ ($\theta=\ang{0}$).
\begin{center}
\begin{tikzpicture}[>=Stealth,scale=1.0]
  \foreach \y in {0.35,0.7}{\draw[blue!60!black,line width=0.8pt,->] (-2.3,\y)--(2.3,\y);}
  \node[blue!70!black,right] at (2.3,0.7){$\vv{B}$};
  \draw[red!80!black,line width=2.2pt] (-2.3,-0.15)--(2.3,-0.15);
  \draw[red!80!black,->,line width=2.2pt] (0.2,-0.15)--(1.1,-0.15); \node[red!80!black,below] at (1.6,-0.15){$I$};
\end{tikzpicture}
\end{center}
Que vaut la force de Laplace subie par le conducteur ? Justifie.
\end{exo}

\begin{exo}[trouver le sens de la force]
Un conducteur horizontal est parcouru par un courant $I$ dirigé \textbf{vers la droite}. Il est plongé
dans un champ $\vv{B}$ dirigé \textbf{vers le haut}, dans le plan de la feuille.
\begin{center}
\begin{tikzpicture}[>=Stealth,scale=1.0]
  \draw[red!80!black,->,line width=2.2pt] (-2,0)--(2,0) node[right]{$I$};
  \draw[blue!70!black,->,line width=1.6pt] (0,-1.2)--(0,1.2) node[above]{$\vv{B}$};
\end{tikzpicture}
\end{center}
En appliquant la règle des trois doigts de la main droite, dans quel sens pointe la force $\vv{F}$ :
sort-elle de la feuille ($\odot$) ou y entre-t-elle ($\otimes$) ?
\end{exo}

\begin{exo}[deux fils, même sens de courant]
Deux fils rectilignes parallèles, dans le plan de la feuille, sont parcourus par des courants de
\textbf{même sens} (tous deux vers le haut).
\begin{center}
\begin{tikzpicture}[>=Stealth,scale=1.0]
  \draw[red!80!black,line width=2pt] (-1,-1.2)--(-1,1.2); \draw[red!80!black,->,line width=2pt] (-1,0)--(-1,0.7); \node[red!80!black,left] at (-1.1,1){$I_1$};
  \draw[red!80!black,line width=2pt] (1,-1.2)--(1,1.2); \draw[red!80!black,->,line width=2pt] (1,0)--(1,0.7); \node[red!80!black,right] at (1.1,1){$I_2$};
\end{tikzpicture}
\end{center}
Les deux fils \textbf{s'attirent}-ils ou se \textbf{repoussent}-ils ?
\end{exo}

\begin{exo}[l'effet de l'angle]
Un conducteur subit une force de Laplace $F_0=\qty{0.2}{\newton}$ lorsqu'il est perpendiculaire au
champ ($\theta=\ang{90}$). Sans changer $B$, $I$ ni $L$ :
\begin{enumerate}[label=\alph*)]
  \item Quelle force subit-il si on l'incline à $\theta=\ang{30}$ ? (On donne $\sin\ang{30}=0{,}5$.)
  \item Quelle force subit-il si on le place parallèle au champ ($\theta=\ang{0}$) ?
\end{enumerate}
\end{exo}

\end{document}
